\section*{第 \ref{ch:g10:reffunc} 章——参考函数}

\begin{solution}{exo:g10:reffunc:1}
$f$：\omterm{def:g10:reffunc:affine}{斜率} $3$，\omterm{def:g10:reffunc:affine}{纵截距} $-2$，\omterm{def:g10:functions:variations}{递增}。\quad
$g$：\omterm{def:g10:reffunc:affine}{斜率} $-\frac12$，\omterm{def:g10:reffunc:affine}{纵截距} $4$，\omterm{def:g10:functions:variations}{递减}。\quad
$h$：\omterm{def:g10:reffunc:affine}{斜率} $0$，\omterm{def:g10:reffunc:affine}{纵截距} $7$，常值\omterm{def:g10:functions:function}{函数}。\quad
$k$：\omterm{def:g10:reffunc:affine}{斜率} $-1$，\omterm{def:g10:reffunc:affine}{纵截距} $0$，\omterm{def:g10:functions:variations}{递减}（且为线性\omterm{def:g10:functions:function}{函数}）。
\end{solution}

\begin{solution}{exo:g10:reffunc:2}
经过 $(2,1)$ 和 $(5,10)$：\omterm{def:g10:reffunc:affine}{斜率} $m = \frac{10 - 1}{5 - 2} = 3$；再
由 $1 = 3 \times 2 + p$ 得 $p = -5$，故 $f(x) = 3x - 5$。

经过 $(-1,4)$ 和 $(3,-4)$：\omterm{def:g10:reffunc:affine}{斜率} $m = \frac{-4 - 4}{3 - (-1)} =
\frac{-8}{4} = -2$；再由 $4 = -2 \times (-1) + p$ 得 $p = 2$：
$f(x) = -2x + 2$。
\end{solution}

\begin{solution}{exo:g10:reffunc:3}
$(3.1)^2 < (3.2)^2$：平方在正数上\omterm{def:g10:functions:variations}{递增}。

$(-2.7)^2 < (-2.8)^2$：平方在负数上\omterm{def:g10:functions:variations}{递减}，
$-2.8 < -2.7$。

$\frac{1}{5.1} > \frac{1}{5.2}$：倒数\omterm{def:g10:functions:function}{函数}在正数上\omterm{def:g10:functions:variations}{递减}。

$\sqrt{17} > \sqrt{15}$：平方根\omterm{def:g10:functions:function}{函数}\omterm{def:g10:functions:variations}{递增}。
\end{solution}

\begin{solution}{exo:g10:reffunc:4}
$x^2 = 16$：有两个解 $x = 4$ 和 $x = -4$（
水平直线 $y = 16$ 与抛物线交于两点）。

$x^2 \leq 16$：抛物线在两交点之间
位于直线下方，即 $x \in \intcc{-4}{4}$。

$x^2 > 9$：抛物线严格高于 $y = 9$ 的部分
在 $\intcc{-3}{3}$ 之外，即 $x \in \intoo{-\infty}{-3} \cup \intoo{3}{+\infty}$。
\end{solution}

\begin{solution}{exo:g10:reffunc:5}
当 $0 < x < 1$ 时，将 $x < 1$ 反复乘以 $x$，
$x^3 < x^2 < x$，由\cref{prop:g10:reffunc:compare}可得 $x < \sqrt x$，
于是
\[
x^3 < x^2 < x < \sqrt x .
\]
取 $x = 0.25$ 检验：$x^3 = 0.015625$，$x^2 = 0.0625$，$x = 0.25$，
$\sqrt x = 0.5$。
\end{solution}

\begin{solution}{exo:g10:reffunc:6}
$\frac1x = 2$ 有唯一解 $x = \frac12$（双曲线
与水平直线 $y = 2$ 在正支上恰有一个交点）。

当 $x > 0$ 时，$\frac1x < 2$。两边乘以 $x > 0$，得 $1 < 2x$，于是
$x > \frac12$，解集为 $\intoo{\frac12}{+\infty}$。

若允许 $x < 0$，则每个负数 $x$ 都满足 $\frac1x < 0 < 2$，
于是完整解集为
$\intoo{-\infty}{0} \cup \intoo{\frac12}{+\infty}$。
\end{solution}

\begin{solution}{exo:g10:reffunc:7}
(a) 在 $\intcc{2}{5}$ 上，所有输入均为正，且平方\omterm{def:g10:functions:variations}{递增}：
$4 \leq x^2 \leq 25$。

(b) 在 $\intcc{-3}{-1}$ 上，平方\omterm{def:g10:functions:variations}{递减}：
最大的平方来自 $-3$：$1 \leq x^2 \leq 9$。

(c) 在 $\intcc{-2}{3}$ 上，\omterm{def:g10:numbers:interval}{区间}包含 $0$，平方在此处
取到\omterm{def:g10:functions:extrema}{最小值} $0$；\omterm{def:g10:functions:extrema}{最大值}是 $(-2)^2 = 4$ 与
$3^2 = 9$ 中的较大者，故 $0 \leq x^2 \leq 9$。
\end{solution}

\begin{solution}{exo:g10:reffunc:8}
第一家公司收费 $f(x) = 1.5x + 4$；第二家公司收费 $g(x) = 2x$。第一家
更便宜的条件是
\[
1.5x + 4 < 2x
\ \Longleftrightarrow\
4 < 0.5x
\ \Longleftrightarrow\
x > 8 .
\]
超过 $8$ 公里时，第一家公司更便宜；恰好 $8$ 公里时，两家收费
均为 $16$。
\end{solution}

\begin{solution}{exo:g10:reffunc:9}
两边均为非负数，因此原不等式等价于
两边平方后所得的不等式：
\[
\left(\sqrt{a+b}\right)^2 = a + b
\qquad\text{且}\qquad
\left(\sqrt a + \sqrt b\right)^2 = a + 2\sqrt a \sqrt b + b .
\]
由于 $2\sqrt a\sqrt b \geq 0$，第二个平方不小于第一个，
从而证明 $\sqrt{a+b} \leq \sqrt a + \sqrt b$。等号成立
当且仅当 $\sqrt a \sqrt b = 0$，即 $a = 0$ 或 $b = 0$。
\end{solution}

\begin{solution}{exo:g10:reffunc:10}
\emph{1.} 将右边通分：
\[
2 + \frac{3}{x-1} = \frac{2(x-1) + 3}{x - 1} = \frac{2x + 1}{x - 1}
= f(x).
\]

\emph{2.} 在 $\intoo{1}{+\infty}$ 上，随 $x$ 增大，$x - 1$ 增大，
且保持为正，因此 $\frac{3}{x-1}$ \omterm{def:g10:functions:variations}{递减}（倒数\omterm{def:g10:functions:function}{函数}乘以
正常数 $3$），故 $f(x) = 2 + \frac{3}{x-1}$ \omterm{def:g10:functions:variations}{递减}：
$f$ 在 $\intoo{1}{+\infty}$ 上严格\omterm{def:g10:functions:variations}{递减}。
\end{solution}

\begin{solution}{pb:g10:reffunc:1}
\textbf{1.} $d(30) = 3^2 = 9$ 米；$d(50) = 25$ 米；$d(90) =
81$ 米；$d(130) = 169$ 米。

\textbf{2.} $v$ 加倍，$\frac v{10}$ 随之加倍，其平方乘以
$4$：$d(100) = 100$ 米 $= 4 \times d(50)$。
速度加倍，距离变为四倍——
正是平方\omterm{def:g10:functions:function}{函数}的伸缩规律。

\textbf{3.} 由 $\left(\frac{v}{10}\right)^2 = 50$ 可得
$\frac{v}{10} = \sqrt{50}$，故 $v = 10\sqrt{50} \approx 71$
千米/时。开方结果唯一，因速度为正，
且在
$\intco{0}{+\infty}$ 上，平方\omterm{def:g10:functions:function}{函数}严格\omterm{def:g10:functions:variations}{递增}
（\cref{prop:g10:reffunc:sqrt}：唯一的正\omterm{def:g10:functions:function}{原像}）。

\textbf{4.} $d(31) - d(30) = 9.61 - 9 = 0.61$ 米，而
$d(91) - d(90) = 82.81 - 81 = 1.81$ 米：同样增加 $+1$ 千米/时，
在高速时所增加的距离约为三倍。抛物线不仅上升，
而且\emph{变陡}——其上升速率随 $x$ 增大。

\textbf{5.} $D(50) = 14 + 25 = 39$ 米；$D(130) = 36.4 + 169 =
205.4$ 米。市区中一次项（反应距离）占比较大，
而在高速公路上，平方项将其压倒——一次项
稳步增长，平方项则迅速拉开距离。

\textbf{6.} $F(0.5) = 240$，$F(1) = 120$，$F(3) = 40$。

\textbf{7.} 倒数\omterm{def:g10:functions:function}{函数}\omterm{def:g10:functions:variations}{递减}：离支点越远，
所需的力越小——若力臂足够长（$d$ 极大），
任何多么微小的力，都能平衡任何
载荷：阿基米德的豪言藏于双曲线的尾部，它
无限趋近于 $0$ 却永不触及。而 $d = 0$ 是
禁止取值：没有力臂，便没有杠杆——即铅直渐近线。

\textbf{8.} $\frac{100}{4} = 25$（距离 $2$ 米处），
$\frac{100}{9} \approx 11$（$3$ 米处），$1$（$10$ 米处）。
距离 $d$ 处，光已铺展于面积为 $4\pi d^2$ 的球面上
（阿基米德墓碑问题，见初中分册）。同一能量被分摊到
随 $d^2$ 增大的面积上——故为\emph{平方反比}关系。

\textbf{9.} 每人分摊 $60$、$30$、$20$ 欧元，对应 $n = 10, 20,
30$。由 $\frac{600}{n} \leq 25$ 得 $n \geq 24$ 人。双曲线承诺，
随 $n$ 增大，人均费用越来越低——
却永远不会降为 $0$：这趟车从来不是免费的。

\textbf{10.} $a(x) = \frac{2x + 3}{x} = 2 + \frac3x$：
在 $\intoo{0}{+\infty}$ 上\omterm{def:g10:functions:variations}{递减}（由倒数\omterm{def:g10:functions:function}{函数}平移而来），
趋近水平渐近线 $y = 2$。从经济角度看：将
固定的 $3$ 欧元摊到更多海报上，使单位成本降低，
逼近材料费中不可压缩的 $2$ 欧元——
这便是规模经济，且存在一个硬性下限。

\textbf{11.} 地球：$a^3 = 1$，$T^2 = 1$。火星：
$1.52^3 \approx 3.51$，$1.88^2 \approx 3.53$。木星：
$5.20^3 \approx 140.6$，$11.86^2 \approx 140.7$。开普勒的
发现：对每颗行星都有 $T^2 = a^3$——
整个天空遵循同一条定律。

\textbf{12.} $T = \sqrt{a^3}$：对土星，
$\sqrt{9.54^3} = \sqrt{868} \approx 29.5$ 年，与观测值
$29.4$ 年相比，定律精确到十分位成立。

\textbf{13.} 小行星：$T = \sqrt{4^3} = \sqrt{64} = 8$ 年。
彗星：$a^3 = 27^2 = 729 = 9^3$，故 $a = 9$ 天文单位。

\textbf{14.} 平方（$T^2$）、立方（$a^3$）与开平方（用来求 $T$）。
伸缩规律：$a$ 乘以 $4$，则 $T$ 乘以
$4\sqrt4 = 8$——小行星印证了这一点（$a$ 为地球的
四倍，$T$ 为八倍）。

\textbf{15.} 一张数表与几个标准\omterm{def:g10:functions:function}{函数}的图形对照，
竟得出宇宙的一条定律，比人们明白\emph{其中原因}
早了七十年——
认出\omterm{def:g10:functions:function}{函数}，便是科学的一半。

\textbf{16.} 在 $\intoo{0}{1}$ 上：$x^2 < x < \sqrt x$；在
$\intoo{1}{+\infty}$ 上：$\sqrt x < x < x^2$。验证：
当 $x = 0.25$ 时：$0.0625 < 0.25 < 0.5$；当 $x = 4$ 时：
$2 < 4 < 16$。

\textbf{17.} 对 $x \geq 0$，两边均 $\geq 0$，故
$\sqrt x > x \iff x > x^2 \iff x(1 - x) > 0 \iff
x \in \intoo{0}{1}$.

\textbf{18.} 当 $x = 0.5$ 时：$x^3 = 0.125 < x^2 = 0.25 < x =
0.5$。当 $x = 2$ 时：$2 < 4 < 8$，大小关系反转。
$x^2$ 与 $x^3$ 的交点：$x^3 - x^2 = x^2(x - 1) = 0$，
即 $x = 0$ 与 $x = 1$ 处。

\textbf{19.} 由 $10t^2 = 100\sqrt t$ 得 $t^2 = 10\sqrt t$；
平方得 $t^4 = 100t$，故 $t^3 = 100$，
$t = \sqrt[3]{100} \approx 4.64$ 年，约 $4$ 年
又 $8$ 个月。（此时两种方案均支付约 $215$ 欧元。）寓意：
根式起步快，幂次终能胜出。

\textbf{20.} 刹车：抛物线的\emph{变陡}，
正是速度的危险。杠杆：\emph{铅直渐近线}位于 $0$ 处，
而长长的尾部，正是技师的优势。
灯：$d^2$ 来自\emph{扩散球面}，决定了光的衰减。开普勒：
太阳系的节拍由\emph{指数} $\frac32$ 决定。
储蓄方案 A：根式的\emph{变平缓}，正是慢速储蓄者的
命运。五张\omterm{def:g10:functions:graph}{图像}，五条定律：自然的字母表，果然如此。
\end{solution}
